\section{2015 Analysis \& Differential Equations}

\subsection{Individual}

\begin{exercise}
\label{analysis:2015:1}
Let $f_n\in L^2(\mathbb{R})$ be a sequence of measurable functions over the line, $f_n\to f$ almost everywhere. Let $\|f_n\|_{L^2}\to\|f\|_{L^2}$. Prove that $\|f_n-f\|_{L^2}\to 0$.
\end{exercise}

\begin{solution}
\textbf{Prerequisites:}
\begin{enumerate}
    \item \textbf{Fatou's Lemma:} For any sequence of non-negative measurable functions $g_n$, the integral of the limit inferior is less than or equal to the limit inferior of the integrals: $\int \liminf_{n\to\infty} g_n(x) \, dx \leq \liminf_{n\to\infty} \int g_n(x) \, dx$. This lemma is a fundamental tool in Lebesgue integration theory for dealing with limits of integrals.
    \item \textbf{Triangle Inequality in $\mathbb{C}$:} For any complex numbers $a$ and $b$, the inequality $|a - b|^2 \leq (|a| + |b|)^2 \leq 2|a|^2 + 2|b|^2$ holds. This provides a bound for the difference of two functions in terms of their magnitudes.
    \item \textbf{$L^2$ Norm Definition:} The $L^2$ norm of a function $f$ on $\mathbb{R}$ is defined by $\|f\|_{L^2} = \left(\int_{\mathbb{R}} |f(x)|^2 dx\right)^{1/2}$. Convergence in $L^2$ means that the "distance" between functions, measured by the integral of the square of their difference, goes to zero.
\end{enumerate}

\textbf{Steps:}
\begin{enumerate}
    \item \textbf{Construct a non-negative helper sequence:} We define a new sequence of functions $g_n(x) = 2(|f_n(x)|^2 + |f(x)|^2) - |f_n(x) - f(x)|^2$. 
    From the triangle inequality bound $|f_n - f|^2 \leq 2|f_n|^2 + 2|f|^2$, we see that $g_n(x) \geq 0$ for all $x$. This non-negativity is the critical requirement for applying Fatou's Lemma in the next step.

    \item \textbf{Determine the pointwise limit:} We are given that $f_n \to f$ almost everywhere (a.e.). This implies that $|f_n(x)|^2 \to |f(x)|^2$ a.e. and $|f_n(x) - f(x)|^2 \to 0$ a.e. 
    Consequently, the pointwise limit inferior of our sequence $g_n$ is:
    \[ \liminf_{n\to\infty} g_n(x) = 2(|f(x)|^2 + |f(x)|^2) - 0 = 4|f(x)|^2 \text{ a.e.} \]

    \item \textbf{Apply Fatou's Lemma:} Integrating the inequality $\int \liminf g_n \leq \liminf \int g_n$, we obtain:
    \[ \int_{\mathbb{R}} 4|f(x)|^2 \, dx \leq \liminf_{n\to\infty} \int_{\mathbb{R}} [2|f_n(x)|^2 + 2|f(x)|^2 - |f_n(x) - f(x)|^2] \, dx \]
    Substituting the norm notation $\|f\|_{L^2}^2 = \int |f|^2$:
    \[ 4\|f\|_{L^2}^2 \leq \liminf_{n\to\infty} [2\|f_n\|_{L^2}^2 + 2\|f\|_{L^2}^2 - \|f_n - f\|_{L^2}^2] \]

    \item \textbf{Utilize the convergence of norms:} We are given that $\|f_n\|_{L^2} \to \|f\|_{L^2}$. As $n \to \infty$, the term $2\|f_n\|_{L^2}^2$ approaches $2\|f\|_{L^2}^2$. 
    The expression inside the limit inferior becomes $4\|f\|_{L^2}^2 - \|f_n - f\|_{L^2}^2$. 
    Taking the limit:
    \[ 4\|f\|_{L^2}^2 \leq 4\|f\|_{L^2}^2 - \limsup_{n\to\infty} \|f_n - f\|_{L^2}^2 \]
    (Note: $\liminf (A - B_n) = A - \limsup B_n$ for a constant $A$).

    \item \textbf{Final Deduction:} Subtracting $4\|f\|_{L^2}^2$ from both sides, we get:
    \[ 0 \leq -\limsup_{n\to\infty} \|f_n - f\|_{L^2}^2 \implies \limsup_{n\to\infty} \|f_n - f\|_{L^2}^2 \leq 0 \]
    Since the norm is always non-negative, the limit must be exactly 0. This proves that $\|f_n - f\|_{L^2} \to 0$.
\end{enumerate}
\par\medskip
\textbf{Intuition:}
\par\medskip
A common pitfall in analysis is assuming that pointwise convergence ($f_n \to f$) implies norm convergence ($\|f_n - f\| \to 0$). This isn't always true; for example, a sequence of narrow "spikes" can stay height 1 but become so thin that they disappear pointwise, yet their integral remains constant.
\par\medskip
However, if the "total energy" (the norm) of the sequence also converges to the energy of the limit function, it prevents energy from "escaping" into these disappearing spikes or moving off to infinity. The $L^2$ norm ensures that if the shapes match and the total energy matches, the functions themselves must be coming together in the mean-square sense.
\end{solution}

\begin{exercise}
\label{analysis:2015:2}
Let $f$ be a continuous function on $[a,b]$, define $M_n = \int_a^b f(x)x^n\,dx$. Suppose that $M_n=0$ for all integers $n\geq 0$. Show that $f(x)=0$ for all $x$.
\end{exercise}

\begin{solution}
\textbf{Prerequisites:}
\begin{enumerate}
    \item \textbf{Weierstrass Approximation Theorem:} This theorem states that any continuous function $f$ on a closed interval $[a,b]$ can be uniformly approximated by polynomials to any degree of accuracy. For any $\epsilon > 0$, there exists a polynomial $P(x)$ such that $|f(x) - P(x)| < \epsilon$ for all $x \in [a,b]$.
    \item \textbf{Linearity of Integration:} The integral of a sum is the sum of the integrals, and constants can be factored out. If a function is orthogonal to $1, x, x^2, \dots$, it is orthogonal to any polynomial (which is just a finite linear combination of these powers).
    \item \textbf{Inner Product Interpretation:} The integral $\int f(x)g(x) dx$ can be thought of as an "inner product" $\langle f, g \rangle$, similar to the dot product in $\mathbb{R}^n$. If $\langle f, f \rangle = \int f^2 = 0$ for a continuous function $f$, then $f$ must be zero everywhere.
\end{enumerate}

\textbf{Steps:}
\begin{enumerate}
    \item \textbf{Orthogonality to all polynomials:} We are given that $\int_a^b f(x)x^n \, dx = 0$ for all $n \geq 0$. By linearity, for any polynomial $P(x) = \sum_{k=0}^N c_k x^k$, we have:
    \[ \int_a^b f(x)P(x) \, dx = \sum_{k=0}^N c_k \int_a^b f(x)x^k \, dx = 0 \]
    This means $f$ is "perpendicular" to the entire space of polynomials.

    \item \textbf{Uniformly approximate $f$:} Since $f$ is continuous on $[a,b]$, we can apply the Weierstrass Approximation Theorem. For a given $\epsilon > 0$, there exists a polynomial $P_\epsilon(x)$ such that:
    \[ \sup_{x\in[a,b]} |f(x) - P_\epsilon(x)| < \epsilon \]

    \item \textbf{Expand the integral of $f^2$:} We want to show $\int f^2 = 0$. We write:
    \[ \int_a^b f(x)^2 \, dx = \int_a^b f(x)(f(x) - P_\epsilon(x) + P_\epsilon(x)) \, dx = \int_a^b f(x)(f(x) - P_\epsilon(x)) \, dx + \int_a^b f(x)P_\epsilon(x) \, dx \]
    By step 1, the second integral $\int f P_\epsilon$ is exactly zero because $P_\epsilon$ is a polynomial.

    \item \textbf{Estimate the remaining term:} Using the uniform bound from step 2:
    \[ \int_a^b f(x)^2 \, dx = \int_a^b f(x)(f(x) - P_\epsilon(x)) \, dx \leq \int_a^b |f(x)| \cdot |f(x) - P_\epsilon(x)| \, dx \]
    \[ \int_a^b f(x)^2 \, dx \leq \int_a^b |f(x)| \cdot \epsilon \, dx = \epsilon \int_a^b |f(x)| \, dx \]

    \item \textbf{Final Conclusion:} Since $\epsilon$ can be chosen arbitrarily small, the non-negative number $\int f^2$ must be smaller than any positive value. Thus, $\int_a^b f(x)^2 \, dx = 0$. Because $f^2$ is a continuous, non-negative function, its integral being zero implies $f(x)^2 = 0$, and thus $f(x) = 0$ for all $x \in [a,b]$.
\end{enumerate}
\par\medskip
\textbf{Intuition:}
\par\medskip
Imagine the space of all continuous functions as a giant room. The polynomials $\{1, x, x^2, \dots\}$ are like a set of "basis vectors" that are so dense they can eventually reach any point in the room. 
\par\medskip
If you find a vector $f$ that is perpendicular to every single one of these basis vectors, it must be the case that $f$ is the zero vector. If $f$ had any "length" in any direction, it would eventually have a non-zero projection onto one of the polynomials. Since it doesn't, it must be zero.
\end{solution}

\begin{exercise}
\label{analysis:2015:3}
Determine all entire functions $f$ that satisfy the inequality
\[
|f(z)| \leq |z|^2|\operatorname{Im}(z)|^2
\]
for $z$ sufficiently large.
\end{exercise}

\begin{solution}
\textbf{Prerequisites:}
\begin{enumerate}
    \item \textbf{Generalized Liouville's Theorem:} A standard version of Liouville's Theorem states that a bounded entire function is constant. The generalized version states that if an entire function $f$ grows no faster than a polynomial of degree $n$ (i.e., $|f(z)| \leq C|z|^n$ for large $|z|$), then $f$ must itself be a polynomial of degree at most $n$.
    \item \textbf{Identity Theorem:} This theorem is a cornerstone of complex analysis. It states that if two holomorphic functions are equal on a set with an accumulation point (like a line segment or a disk), they are equal everywhere on their domain. For a polynomial, this simply means that if it has more roots than its degree, it must be the zero polynomial.
    \item \textbf{Imaginary Part Relationship:} For any complex number $z = x + iy$, the imaginary part is bounded by the magnitude: $|\operatorname{Im}(z)| = |y| \leq \sqrt{x^2 + y^2} = |z|$.
\end{enumerate}

\textbf{Steps:}
\begin{enumerate}
    \item \textbf{Establish a growth bound:} The problem gives us $|f(z)| \leq |z|^2 |\operatorname{Im}(z)|^2$. Using the fact that $|\operatorname{Im}(z)| \leq |z|$, we can see that:
    \[ |f(z)| \leq |z|^2 \cdot |z|^2 = |z|^4 \]
    for sufficiently large $z$. By the Generalized Liouville's Theorem, $f(z)$ must be a polynomial of degree at most 4.

    \item \textbf{Analyze the function on the real axis:} Let $z = x$ where $x$ is a real number. For any real number, the imaginary part is zero: $\operatorname{Im}(x) = 0$.
    Plugging this into the original inequality:
    \[ |f(x)| \leq |x|^2 \cdot |0|^2 = 0 \]
    This means that $f(x) = 0$ for all real $x$ that are sufficiently large.

    \item \textbf{Apply the Identity Theorem:} We have discovered that $f(z)$ is a polynomial that is equal to zero on an infinite set of points (the large real numbers). 
    A polynomial of degree at most 4 can have at most 4 roots unless it is the zero polynomial. Since our $f(z)$ has infinitely many roots on the real line, we must conclude that $f(z) = 0$ for all $z \in \mathbb{C}$.

    \item \textbf{Verification:} If $f(z) = 0$, then $|0| \leq |z|^2 |\operatorname{Im}(z)|^2$ is $0 \leq |z|^2 |\operatorname{Im}(z)|^2$, which is clearly true for all $z$. Thus, the zero function is the only solution.
\end{enumerate}
\par\medskip
\textbf{Intuition:}
\par\medskip
Holomorphic functions are "rigid." Unlike smooth real functions, you cannot just "squash" a part of a holomorphic function to zero without affecting the rest of it. 
\par\medskip
The growth condition $|f(z)| \leq |z|^2 |\operatorname{Im}(z)|^2$ forces the function to be zero whenever the imaginary part is zero. Since the imaginary part is zero on the entire real line, the function is forced to be zero on a "large" part of the complex plane. This "squashing" on the real line, combined with the limited growth at infinity, leaves the function with no choice but to be zero everywhere.
\end{solution}


\begin{exercise}
\label{analysis:2015:4}
Describe all functions that are holomorphic over the unit disk $D=\{|z|<1\}$, continuous on $\bar{D}$, and map the boundary of the disk into the boundary of the disk.
\end{exercise}

\begin{solution}
\textbf{Prerequisites:}
\begin{enumerate}
    \item \textbf{Maximum Modulus Principle:} If a function is holomorphic in a domain and continuous on its boundary, its maximum absolute value is achieved on the boundary. If the maximum is achieved in the interior, the function must be constant.
    \item \textbf{Blaschke Factors:} A function of the form $B_a(z) = \frac{z - a}{1 - \bar{a} z}$ for some $|a| < 1$ is a holomorphic map from the unit disk $D$ to itself. Crucially, on the boundary ($|z|=1$), we have $|B_a(z)| = 1$. These factors are the building blocks for functions mapping the circle to itself.
    \item \textbf{Minimum Modulus Principle:} Similar to the maximum principle, if a holomorphic function has no zeros in a domain, its minimum absolute value must be achieved on the boundary.
\end{enumerate}
\par\medskip
\textbf{Steps:}
\begin{enumerate}
    \item \textbf{Bound the function:} We are given that $|f(z)| = 1$ whenever $|z| = 1$. By the Maximum Modulus Principle, we must have $|f(z)| \leq 1$ for all $z$ inside the disk $D$.
    \item \textbf{Check for zeros:} If $f(z)$ has no zeros inside the disk, then the function $g(z) = 1/f(z)$ is also holomorphic in the disk and continuous on the boundary. Since $|g(z)| = 1/1 = 1$ on the boundary, the Maximum Modulus Principle applied to $g$ tells us $|g(z)| \leq 1$ inside, which means $|f(z)| \geq 1$. Combined with step 1, this forces $|f(z)| = 1$ everywhere, making $f$ a constant of modulus 1 ($f(z) = e^{i\theta}$).
    \item \textbf{Identify the zeros:} If $f(z)$ is not constant, it must have at least one zero in the disk. Since $f$ is continuous on the closed disk and $|f|=1$ on the boundary, it can only have a finite number of zeros, say $a_1, a_2, \dots, a_n$ (counted with multiplicity). If it had infinitely many, they would accumulate in the disk or on the boundary, which would force $f=0$.
    \item \textbf{Factor out the zeros:} Define a Blaschke product using these zeros:
    \[ B(z) = \prod_{k=1}^n \frac{z - a_k}{1 - \bar{a}_k z} \]
    Consider the function $h(z) = f(z) / B(z)$. This function is holomorphic in $D$ (the zeros of $B$ cancel the zeros of $f$) and has no zeros in $D$.
    \item \textbf{Determine the final form:} On the boundary $|z|=1$, we have $|B(z)| = 1$, so $|h(z)| = |f(z)|/|B(z)| = 1/1 = 1$. Applying the logic from step 2 to the non-vanishing function $h(z)$, we conclude that $h(z)$ must be a constant of modulus 1. Therefore, $h(z) = e^{i\theta}$.
    \item \textbf{Conclusion:} All such functions are of the form:
    \[ f(z) = e^{i\theta} \prod_{k=1}^n \frac{z - a_k}{1 - \bar{a}_k z} \]
    where $n \geq 0$, $\theta \in \mathbb{R}$, and $|a_k| < 1$. These are exactly the finite Blaschke products.
\end{enumerate}
\par\medskip
\textbf{Intuition:}
\par\medskip
Think of the unit disk as a piece of elastic fabric. The problem asks for all ways to transform this fabric holomorphically such that the edge of the fabric stays exactly on the boundary circle. 
\par\medskip
A single Blaschke factor $\frac{z-a}{1-\bar{a}z}$ essentially takes the disk and "shifts" its center to $a$ while keeping the boundary fixed. Multiplying these factors together "wraps" the disk around itself multiple times. The number of zeros $n$ tells you how many times the function "covers" the disk. The requirement of holomorphicity is so strong that these "wrapping" maps are the only possible solutions.
\end{solution}

\begin{exercise}
\label{analysis:2015:5}
Let $T:H_1\to H_2$ and $Q:H_2\to H_1$ be bounded linear operators of Hilbert spaces $H_1,H_2$. Let
\[
QT=I-S_1,\qquad TQ=I-S_2,
\]
where \(S_1\) and \(S_2\) are compact operators. Prove that
\[
\ker T=\{v\in H_1:Tv=0\},\qquad
\operatorname{coker}T=H_2/\operatorname{Im}T
\]
are finite-dimensional and that
\[
\operatorname{Im}T=\{Tv:v\in H_1\}
\]
is closed in \(H_2\).
\end{exercise}

\begin{solution}
\textbf{Finite dimensionality of the kernel.}
If \(v\in\ker T\), then
\[
v=QTv+S_1v=S_1v.
\]
Thus \(S_1\) is the identity on \(\ker T\). If \(\ker T\) were infinite-dimensional, its unit ball would not be compact, but its image under the compact operator \(S_1\) would be relatively compact. This contradiction shows
\[
\dim\ker T<\infty.
\]

\textbf{Closedness of the range.}
Suppose \(Tv_n\to w\) in \(H_2\). Decompose
\[
v_n=k_n+u_n,\qquad k_n\in\ker T,\quad u_n\in(\ker T)^\perp.
\]
Since \(Tk_n=0\), it is enough to work with \(u_n\). We claim that \(\{u_n\}\) is bounded. If not, set \(y_n=u_n/\|u_n\|\) along a subsequence with \(\|u_n\|\to\infty\). Then \(\|y_n\|=1\), \(y_n\in(\ker T)^\perp\), and
\[
Ty_n=\frac{Tu_n}{\|u_n\|}\to0.
\]
Using \(QT=I-S_1\),
\[
y_n=QTy_n+S_1y_n.
\]
Compactness of \(S_1\) gives a convergent subsequence of \(S_1y_n\), hence a convergent subsequence \(y_{n_j}\to y\). Passing to the limit gives \(Ty=0\), so \(y\in\ker T\). But \(y\in(\ker T)^\perp\) and \(\|y\|=1\), a contradiction. Therefore \(\{u_n\}\) is bounded.

Now \(Tu_n=Tv_n\to w\), and
\[
u_n=QTu_n+S_1u_n.
\]
After passing to a subsequence, \(S_1u_n\) converges by compactness, while \(QTu_n\to Qw\). Hence \(u_n\to u\) for some \(u\in H_1\). By continuity,
\[
Tu=w.
\]
Thus \(w\in\operatorname{Im}T\), so \(\operatorname{Im}T\) is closed.

\textbf{Finite dimensionality of the cokernel.}
For a bounded operator with closed range,
\[
\operatorname{coker}T\cong(\operatorname{Im}T)^\perp=\ker T^*.
\]
Taking adjoints in the parametrix identities gives
\[
T^*Q^*=I-S_1^*,\qquad Q^*T^*=I-S_2^*,
\]
and \(S_1^*,S_2^*\) are compact. Applying the kernel argument to \(T^*\) gives
\[
\dim\ker T^*<\infty.
\]
Hence \(\operatorname{coker}T\) is finite-dimensional.
\end{solution}


\begin{exercise}
\label{analysis:2015:6}
Let $H_1$ be the Sobolev space on the unit interval $[0,1]$, i.e., the Hilbert space consisting of functions $f\in L^2([0,1])$ such that
\[
\|f\|_1^2 = \sum_{n=-\infty}^{\infty}(1+n^2)|\hat{f}(n)|^2 < \infty,
\]
where
\[
\hat{f}(n) = \frac{1}{2\pi}\int_0^1 f(x)e^{-2\pi inx}\,dx
\]
are Fourier coefficients of $f$. Show that there exists constant $C>0$ such that
\[
\|f\|_{L^{\infty}} \leq C\|f\|_1
\]
for all $f\in H_1$, where $\|.\|_{L^{\infty}}$ stands for the usual supremum norm. (Hint: Use Fourier series.)
\end{exercise}

\begin{solution}
\textbf{Prerequisites:}
\begin{enumerate}
    \item \textbf{Fourier Series Expansion:} For any function in $L^2([0,1])$, we can represent it as a sum of complex exponentials: $f(x) = \sum_{n=-\infty}^\infty \hat{f}(n) e^{2\pi i n x}$. This is the "frequency domain" representation of the function.
    \item \textbf{Cauchy-Schwarz Inequality for Series:} This is a fundamental inequality stating that for any two sequences $a_n$ and $b_n$, we have $\left(\sum |a_n b_n|\right)^2 \leq \left(\sum |a_n|^2\right) \left(\sum |b_n|^2\right)$. 
    \item \textbf{Sobolev Norm:} The norm defined by $\|f\|_1^2 = \sum (1+n^2)|\hat{f}(n)|^2$ is called the $H^1$ Sobolev norm. It measures both the average magnitude of the function (the $L^2$ part) and its "smoothness" (since the $n^2$ factor heavily weights the high-frequency components that represent rapid oscillations or jumps).
\end{enumerate}

\textbf{Steps:}
\begin{enumerate}
    \item \textbf{Express $|f(x)|$ in terms of Fourier coefficients:} 
    We know that $f(x) = \sum_{n=-\infty}^\infty \hat{f}(n) e^{2\pi i n x}$. Applying the triangle inequality to this sum:
    \[ |f(x)| = \left| \sum_{n=-\infty}^\infty \hat{f}(n) e^{2\pi i n x} \right| \leq \sum_{n=-\infty}^\infty |\hat{f}(n)| \cdot |e^{2\pi i n x}| \]
    Since $|e^{i\theta}| = 1$ for any real $\theta$, we have $|f(x)| \leq \sum_{n=-\infty}^\infty |\hat{f}(n)|$.

    \item \textbf{Apply the Cauchy-Schwarz trick:} We want to introduce the factor $\sqrt{1+n^2}$ to relate this sum to the given Sobolev norm. We multiply and divide each term by $\sqrt{1+n^2}$:
    \[ \sum_{n=-\infty}^\infty |\hat{f}(n)| = \sum_{n=-\infty}^\infty \left( |\hat{f}(n)| \sqrt{1+n^2} \cdot \frac{1}{\sqrt{1+n^2}} \right) \]
    Using Cauchy-Schwarz with $a_n = |\hat{f}(n)| \sqrt{1+n^2}$ and $b_n = \frac{1}{\sqrt{1+n^2}}$, we get:
    \[ \sum_{n=-\infty}^\infty |\hat{f}(n)| \leq \sqrt{\sum_{n=-\infty}^\infty (1+n^2)|\hat{f}(n)|^2} \cdot \sqrt{\sum_{n=-\infty}^\infty \frac{1}{1+n^2}} \]

    \item \textbf{Identify the constants:}
    The first square root is precisely the definition of the Sobolev norm $\|f\|_1$. 
    The second square root is a constant $C = \sqrt{\sum_{n=-\infty}^\infty \frac{1}{1+n^2}}$. This series converges because the terms are approximately $1/n^2$, and we know the sum of reciprocals of squares converges.

    \item \textbf{Conclude:}
    We have shown that for any $x$, $|f(x)| \leq C \|f\|_1$. 
    Since this bound is independent of $x$, it holds for the supremum of $f$ as well: 
    $\|f\|_{L^\infty} = \sup_{x} |f(x)| \leq C \|f\|_1$.
\end{enumerate}
\par\medskip
\textbf{Intuition:}
\par\medskip
Why does knowing a "smoothness" norm give you a bound on the "maximum height" of a function? 
\par\medskip
Think about a function that has a very high, narrow spike (like a Dirac delta). Such a spike contains a huge amount of high-frequency energy (its Fourier coefficients $\hat{f}(n)$ do not decay). 
\par\medskip
The Sobolev norm $\|f\|_1$ "punishes" high frequencies by multiplying them by $n^2$. If $\|f\|_1$ is finite, it means the high-frequency components must die out very quickly. If the high frequencies die out quickly, the function is prevented from forming sharp spikes, thus keeping the overall height $\|f\|_{L^\infty}$ under control. This is a simple version of the Sobolev Embedding Theorem.
\end{solution}


\subsection{Team}

\begin{exercise}
\label{analysis:2015:7}
Let $\phi\in C([a,b],\mathbb{R})$. Suppose for every function $h\in C^1([a,b],\mathbb{R})$, $h(a)=h(b)=0$, we have
\[
\int_a^b \phi(x)h(x)\,dx = 0.
\]
Prove that $\phi(x)=0$.
\end{exercise}

\begin{solution}
\textbf{Prerequisites:}
\par\medskip
1. \textbf{Definition of Continuity:} A function $\phi(x)$ is continuous at $x_0$ if for any $\epsilon > 0$, there is a $\delta > 0$ such that $|x - x_0| < \delta \implies |\phi(x) - \phi(x_0)| < \epsilon$. This implies that if a continuous function is non-zero at a point, it must remain non-zero in a small neighborhood around that point.
\par\medskip
2. \textbf{Test Functions (Bump Functions):} In calculus of variations, we often use "test functions" $h(x)$ that are smooth and vanish at the boundaries. A common construction for such a function on an interval $[c, d]$ is $h(x) = (x-c)^2 (x-d)^2$ for $x \in [c, d]$ and $0$ otherwise.
\par\medskip
3. \textbf{Integral of a Positive Function:} If $g(x)$ is a continuous function such that $g(x) \geq 0$ for all $x$ and $g(x_0) > 0$ for at least one point, then $\int g(x) dx > 0$.
\par\medskip
\textbf{Steps:}
\par\medskip
1. \textbf{Assumption for contradiction:} Suppose that $\phi(x)$ is not identically zero on $[a,b]$. This means there exists some point $x_0 \in (a,b)$ such that $\phi(x_0) \neq 0$. Without loss of generality, let us assume $\phi(x_0) > 0$.
\par\medskip
2. \textbf{Identify a positive neighborhood:} Since $\phi$ is continuous, there exists a small distance $\delta > 0$ such that $\phi(x) > \frac{1}{2}\phi(x_0) > 0$ for all $x$ in the interval $[x_0 - \delta, x_0 + \delta] \subset (a,b)$. 
\par\medskip
3. \textbf{Construct a specific test function:} We choose a function $h(x)$ that is strictly positive inside this small interval and zero everywhere else. A suitable choice is:
\[ h(x) = \begin{cases} (x - (x_0-\delta))^2 (x - (x_0+\delta))^2 & \text{if } x \in [x_0 - \delta, x_0 + \delta] \\ 0 & \text{otherwise} \end{cases} \]
Note that $h(x)$ is $C^1$ (its derivative is zero at the transition points) and $h(a) = h(b) = 0$, so it satisfies the conditions for a valid test function.
\par\medskip
4. \textbf{Evaluate the integral:} We compute the integral of the product $\phi(x)h(x)$:
\[ \int_a^b \phi(x)h(x) \, dx = \int_{x_0-\delta}^{x_0+\delta} \phi(x)h(x) \, dx \]
Inside this sub-interval, both $\phi(x)$ and $h(x)$ are strictly positive. Thus, their product is strictly positive, and the integral must be strictly greater than zero.
\par\medskip
5. \textbf{Contradiction and Conclusion:} This result ($\int \phi h > 0$) directly contradicts the given hypothesis that $\int_a^b \phi(x)h(x) \, dx = 0$ for \textit{all} such functions $h$. 
Therefore, our initial assumption that $\phi(x_0) \neq 0$ must be false. We conclude that $\phi(x) = 0$ for all $x \in (a,b)$, and by continuity, $\phi(a) = \phi(b) = 0$ as well.
\par\medskip
\textbf{Intuition:}
\par\medskip
Think of the test function $h(x)$ as a "localized probe." The condition $\int \phi h = 0$ for every possible $h$ means that $\phi$ must be "balanced" or "invisible" to every probe you throw at it. 
\par\medskip
If $\phi$ had any "lump" (a region where it is positive or negative), you could design a probe $h$ that specifically targets that lump (by being positive only there). If you did that, the integral would "catch" the lump and return a non-zero value. The only way $\phi$ can hide from every possible probe is if it has no lumps at all—meaning it must be zero everywhere.
\end{solution}

\begin{exercise}
\label{analysis:2015:8}
Let $f$ be a Lebesgue integrable function over $[a,b+\delta]$, $\delta>0$. Prove that
\[
\lim_{h\to 0^+}\int_a^b |f(x+h)-f(x)|\,dx \to 0.
\]
\end{exercise}

\begin{solution}
\textbf{Prerequisites:}
\par\medskip
1. \textbf{Density of Continuous Functions in $L^1$:} A key result in measure theory states that continuous functions with compact support are "dense" in $L^1$. This means that for any integrable function $f$ and any $\epsilon > 0$, we can find a continuous function $g$ such that $\int |f - g| < \epsilon$.
\par\medskip
2. \textbf{Uniform Continuity:} Any continuous function on a closed, bounded interval $[a, b]$ is uniformly continuous. This means that for any $\epsilon > 0$, there is a $\delta > 0$ such that if $|x - y| < \delta$, then $|g(x) - g(y)| < \epsilon$ for \textit{all} $x, y$ in the interval simultaneously.
\par\medskip
3. \textbf{Triangle Inequality for Integrals:} $\int |A + B + C| \leq \int |A| + \int |B| + \int |C|$. This allows us to break the difference $|f(x+h) - f(x)|$ into manageable pieces by introducing an intermediate function $g$.
\par\medskip
\textbf{Steps:}
\begin{enumerate}
    \item \textbf{Introduce an approximating function:} Let $\epsilon > 0$. Since continuous functions are dense in $L^1$, we can choose a continuous function $g$ on $[a, b+\delta]$ such that the "error" in $L^1$ is small:
    \[ \int_a^{b+\delta} |f(t) - g(t)| \, dt < \frac{\epsilon}{3} \]

    \item \textbf{Decompose the integral:} We use the triangle inequality to insert $g$ into our integral:
    \[ \int_a^b |f(x+h) - f(x)| \, dx \leq \int_a^b |f(x+h) - g(x+h)| \, dx + \int_a^b |g(x+h) - g(x)| \, dx + \int_a^b |g(x) - f(x)| \, dx \]

    \item \textbf{Bound the "approximation" terms:} 
    The third term $\int_a^b |g(x) - f(x)| \, dx$ is less than $\epsilon/3$ by our choice of $g$ in Step 1.
    The first term $\int_a^b |f(x+h) - g(x+h)| \, dx$ is also bounded. By shifting the variable $u = x+h$, we get $\int_{a+h}^{b+h} |f(u) - g(u)| \, du$. For small $h$, this is just a sub-integral of the one in Step 1, so it is also less than $\epsilon/3$.

    \item \textbf{Bound the "continuity" term:} 
    Since $g$ is continuous on the compact set $[a, b+\delta]$, it is uniformly continuous. For our given $\epsilon$, there exists an $H > 0$ such that if $0 < h < H$, then $|g(x+h) - g(x)| < \frac{\epsilon}{3(b-a)}$ for all $x \in [a,b]$.
    Integrating this constant bound over the interval gives:
    \[ \int_a^b |g(x+h) - g(x)| \, dx < \int_a^b \frac{\epsilon}{3(b-a)} \, dx = \frac{\epsilon}{3} \]

    \item \textbf{Combine the bounds:} 
    Adding the three pieces together, we find that for $0 < h < H$:
    \[ \int_a^b |f(x+h) - f(x)| \, dx < \frac{\epsilon}{3} + \frac{\epsilon}{3} + \frac{\epsilon}{3} = \epsilon \]
    Since this holds for any $\epsilon > 0$, the limit as $h \to 0^+$ must be zero.
\end{enumerate}
\par\medskip
\textbf{Intuition:}
\par\medskip
This property is called "continuity in the mean." It says that if you shift an integrable function by a tiny amount, the "average" difference between the original and the shifted version is also tiny. 
\par\medskip
If $f$ were a simple continuous function, this would be obvious because the values themselves wouldn't change much. The power of this proof is that it works even for "messy" functions (like those with many jumps or singularities). We use the fact that even the messiest integrable function can be "blurred" or approximated by a smooth, continuous one. Since the property is clearly true for the smooth function and the error in approximation is small, the property must also be true for the original messy function.
\end{solution}

\begin{exercise}
\label{analysis:2015:9}
Let $L(q,q',t)$ be a function of $(q,q',t)\in TU\times\mathbb{R}$, $U$ is an open domain in $\mathbb{R}^n$. Let $\gamma:[a,b]\to U$ be a curve in $U$. Define a functional
\[
S(\gamma) = \int_a^b L(\gamma(t),\gamma'(t),t)\,dt.
\]
We say that $\gamma$ is an extremal if for every smooth variation of $\gamma$, $\phi(t,s)$, $s\in(-\delta,\delta)$, $\phi(t,0)=\gamma(t)$, $\phi_s=\phi(t,s)$, we have $\frac{dS(\phi_s)}{ds}|_{s=0}=0$. Prove that every extremal $\gamma$ satisfies
\[
\frac{d}{dt}\left(\frac{\partial L}{\partial q'}\right) = \frac{\partial L}{\partial q}.
\]
\end{exercise}

\begin{solution}
\textbf{Prerequisites:}
\par\medskip
1. \textbf{Definition of Variation:} A variation of a curve $\gamma(t)$ is a family of nearby curves $\phi(t,s)$ that all start and end at the same points. The parameter $s$ controls how far we "perturb" the original curve.
\par\medskip
2. \textbf{Integration by Parts:} The formula $\int u v' dt = uv - \int v u' dt$ is crucial for shifting the derivative from one function (the perturbation) to another (the Lagrangian's derivatives).
\par\medskip
3. \textbf{Fundamental Lemma of Calculus of Variations:} As seen in Problem 1 of this section, if $\int_a^b \Phi(t) \eta(t) dt = 0$ for all smooth functions $\eta$ that vanish at the boundaries, then $\Phi(t)$ must be zero.
\par\medskip
\textbf{Steps:}
\begin{enumerate}
    \item \textbf{Define the derivative of the functional:} We are given that $\gamma$ is an extremal, meaning the first derivative of the action $S(\phi_s)$ with respect to the perturbation parameter $s$ is zero at $s=0$. 
    Let $\eta(t) = \frac{\partial \phi}{\partial s}(t,0)$ be the "direction" of the variation. Since the variations fix the endpoints, we must have $\eta(a) = \eta(b) = 0$.

    \item \textbf{Differentiate inside the integral:} We calculate $\frac{d}{ds} S(\phi_s)|_{s=0}$:
    \[ \frac{d}{ds} \int_a^b L(\phi(t,s), \phi'(t,s), t) \, dt \bigg|_{s=0} = \int_a^b \left( \frac{\partial L}{\partial q} \frac{\partial \phi}{\partial s} + \frac{\partial L}{\partial q'} \frac{\partial \phi'}{\partial s} \right) dt = 0 \]
    Plugging in our definition of $\eta(t)$ and noting that $\frac{\partial \phi'}{\partial s} = \frac{\partial}{\partial s} \frac{\partial \phi}{\partial t} = \frac{\partial}{\partial t} \frac{\partial \phi}{\partial s} = \eta'(t)$, we get:
    \[ \int_a^b \left( \frac{\partial L}{\partial q} \eta(t) + \frac{\partial L}{\partial q'} \eta'(t) \right) dt = 0 \]

    \item \textbf{Apply Integration by Parts:} To group everything by $\eta(t)$, we integrate the second term by parts:
    \[ \int_a^b \frac{\partial L}{\partial q'} \eta'(t) \, dt = \left[ \frac{\partial L}{\partial q'} \eta(t) \right]_a^b - \int_a^b \frac{d}{dt} \left( \frac{\partial L}{\partial q'} \right) \eta(t) \, dt \]
    Because $\eta(a) = \eta(b) = 0$, the boundary term $[\dots]_a^b$ vanishes.

    \item \textbf{Combine the terms:} Substituting back into the integral equation:
    \[ \int_a^b \left( \frac{\partial L}{\partial q} - \frac{d}{dt} \left( \frac{\partial L}{\partial q'} \right) \right) \eta(t) \, dt = 0 \]

    \item \textbf{Invoke the Fundamental Lemma:} Since this equation holds for every smooth variation $\eta(t)$ that vanishes at the endpoints, the term inside the parentheses must be identically zero:
    \[ \frac{\partial L}{\partial q} - \frac{d}{dt} \left( \frac{\partial L}{\partial q'} \right) = 0 \implies \frac{d}{dt} \left( \frac{\partial L}{\partial q'} \right) = \frac{\partial L}{\partial q} \]
    This is the celebrated Euler-Lagrange equation.
\end{enumerate}
\par\medskip
\textbf{Intuition:}
\par\medskip
The principle of stationary action is like finding the bottom of a bowl. If you are at the lowest point, a tiny step in any direction shouldn't change your altitude (the action) much—it's "flat" there. 
\par\medskip
By checking that the change in action is zero for \textit{any} tiny perturbation $\eta(t)$, we are effectively checking that the "gradient" of the action is zero. The integration by parts trick is just a way to shift the perspective from "how the perturbation changes" to "how the physics of the curve changes." The resulting Euler-Lagrange equation tells us the exact balance of forces (or rates of change) required for a curve to be the most "efficient" path between two points.
\end{solution}

\begin{exercise}
\label{analysis:2015:10}
Let $f:U\to U$ be a holomorphic function with $U$ a bounded domain in the complex plane. Assuming $0\in U$, $f(0)=0$, $f'(0)=1$, prove that $f(z)=z$.
\end{exercise}

\begin{solution}
\textbf{Prerequisites:}
\par\medskip
1. \textbf{Power Series Expansion:} Any holomorphic function $f(z)$ can be expressed as a power series $f(z) = \sum_{n=0}^\infty a_n z^n$ in a neighborhood of $z=0$.
\par\medskip
2. \textbf{Cauchy's Inequality for Derivatives:} This inequality provides a bound on the derivatives of a holomorphic function: $|f^{(k)}(0)| \leq \frac{k! M}{r^k}$, where $M$ is the maximum value of $|f|$ on a circle of radius $r$ around the origin.
\par\medskip
3. \textbf{Iterates of a Function:} The $n$-th iterate of a function $f$, denoted $f_n$, is the function obtained by applying $f$ repeatedly $n$ times ($f \circ f \circ \dots \circ f$). If $f$ maps a bounded domain $U$ to itself, then all its iterates also map $U$ to itself and are thus uniformly bounded by the same constant $M$.
\par\medskip
\textbf{Steps:}
\begin{enumerate}
    \item \textbf{Assume a non-linear form for $f(z)$:} Suppose for contradiction that $f(z) \neq z$. Since $f(0)=0$ and $f'(0)=1$, the power series expansion of $f(z)$ around the origin must look like:
    \[ f(z) = z + a_k z^k + \dots \]
    where $k \geq 2$ is the index of the first non-zero coefficient after the linear term.

    \item \textbf{Analyze the iterates of $f$:} Consider the $n$-th iterate $f_n(z) = (f \circ f \circ \dots \circ f)(z)$. Let's see what happens to the $z^k$ term when we compose the function. 
    For $n=2$: $f(f(z)) = (z + a_k z^k + \dots) + a_k (z + a_k z^k + \dots)^k + \dots = z + 2a_k z^k + \dots$ 
    By induction, the $n$-th iterate will have the form:
    \[ f_n(z) = z + n a_k z^k + \dots \]

    \item \textbf{Use the boundedness of the domain:} We are given that $f: U \to U$ and $U$ is a bounded domain. This means there exists a constant $M$ such that $|f_n(z)| \leq M$ for all $z \in U$ and for all iterates $n$.

    \item \textbf{Apply Cauchy's Inequality:} Let $D(0, r)$ be a small disk centered at the origin that is contained within $U$. The $k$-th derivative of the $n$-th iterate at $z=0$ is:
    \[ f_n^{(k)}(0) = k! \cdot (n a_k) \]
    Using Cauchy's Inequality:
    \[ |k! \cdot n a_k| \leq \frac{k! M}{r^k} \implies |n a_k| \leq \frac{M}{r^k} \]

    \item \textbf{Obtain a contradiction:} The inequality $|n a_k| \leq \frac{M}{r^k}$ must hold for \textit{every} positive integer $n$. 
    However, if $a_k \neq 0$, then as $n \to \infty$, the left side $|n a_k|$ grows to infinity, while the right side $\frac{M}{r^k}$ remains a fixed constant. This is impossible.

    \item \textbf{Conclusion:} The only way to avoid this contradiction is if $a_k = 0$ for all $k \geq 2$. Thus, the power series must be simply $f(z) = z$.
\end{enumerate}
\par\medskip
\textbf{Intuition:}
\par\medskip
This is a "stretching" argument. If a function has a term like $a_k z^k$, then repeating the function $n$ times "multiplies" that term's influence. 
\par\medskip
Think of $f(z)$ as a map that slightly distorts a bounded shape. If the distortion is non-zero, then repeating that distortion millions of times would eventually stretch the shape so much that it would "break" or fly out of its original bounded container $U$. Since the function is required to keep the shape inside the container forever, the only possible "distortion" is no distortion at all—the identity map.
\end{solution}

\begin{exercise}
\label{analysis:2015:11}
Let $T:H_1\to H_2$ be a bounded operator of Hilbert spaces $H_1,H_2$. Let $S:H_1\to H_2$ be a compact operator, that is, for every bounded sequence $\{v_n\}\in H_1$, $Sv_n$ has a converging subsequence. Show that $\operatorname{coker}(T+S) = H_2/\operatorname{Im}(T+S)$ is finite dimensional and $\operatorname{Im}(T+S)$ is closed in $H_2$. (Hint: Consider equivalent statements in terms of adjoint operators.)
\end{exercise}

\begin{solution}
\textbf{As written, the statement is false.}
Take $H_1=H_2$ to be any infinite-dimensional Hilbert space, and take
\[
T=0,\qquad S=0.
\]
Then $S$ is compact, but
\[
\operatorname{Im}(T+S)=\{0\},\qquad
\operatorname{coker}(T+S)=H_2
\]
is infinite-dimensional. Hence the printed problem is missing a hypothesis on $T$.

\textbf{Correct common version.}
If the intended assumption is that $T$ is a Fredholm operator, then the result is the standard theorem that compact perturbations of Fredholm operators are Fredholm. In the special case implicitly used by the old solution, $H_1=H_2=H$ and $T=I$, the proof is as follows.

Let $A=I+S$.

1. \textbf{Finite dimensionality of $\ker A$.}
If $Av=0$, then $v=-Sv$. On $\ker A$, the compact operator $S$ acts as $-I$. Therefore the unit ball of $\ker A$ is mapped by $S$ to itself up to sign and is relatively compact. By Riesz's compactness criterion, $\ker A$ is finite-dimensional.

2. \textbf{Closedness of $\operatorname{Im}A$.}
Suppose $Av_n\to w$. Decompose
\[
v_n=k_n+u_n,\qquad k_n\in\ker A,\quad u_n\in(\ker A)^\perp.
\]
Since $Ak_n=0$, it is enough to work with $u_n$. If $\{u_n\}$ were unbounded, set $y_n=u_n/\|u_n\|$. Then $\|y_n\|=1$, $y_n\in(\ker A)^\perp$, and
\[
Ay_n=\frac{Au_n}{\|u_n\|}\to 0.
\]
Compactness of $S$ gives a subsequence with $Sy_{n_j}$ convergent. Since
\[
y_{n_j}=Ay_{n_j}-Sy_{n_j},
\]
the sequence $y_{n_j}$ converges to some $y$ with $\|y\|=1$. Passing to the limit in $Ay_{n_j}\to0$ gives $Ay=0$, so $y\in\ker A$. But also $y\in(\ker A)^\perp$, a contradiction. Thus $\{u_n\}$ is bounded.

Now compactness of $S$ gives a subsequence $Su_{n_j}\to z$, and
\[
u_{n_j}=Au_{n_j}-Su_{n_j}\to w-z=:u.
\]
By continuity, $Au=w$, so $w\in\operatorname{Im}A$. Therefore the range is closed.

3. \textbf{Finite dimensionality of the cokernel.}
For an operator with closed range,
\[
\operatorname{coker}A\cong(\operatorname{Im}A)^\perp=\ker A^*.
\]
Since $A^*=I+S^*$ and $S^*$ is compact, the first part applied to $A^*$ gives $\dim\ker A^*<\infty$. Hence $\operatorname{coker}A$ is finite-dimensional.

The same argument, after reducing by a parametrix, proves the general Fredholm case.
\end{solution}

\begin{exercise}
\label{analysis:2015:12}
Let $u\in C^2(\Omega)$, where $\Omega\subset\mathbb{R}^d$ is a bounded domain with a smooth boundary.

\begin{itemize}
\item[a)] Let $u$ be a solution of the equation $\Delta u = f$, $u|_{\partial\Omega}=0$, $f\in L^2(\Omega)$. Prove that there is a constant $C$ depending only on $\Omega$ such that
\[
\int_{\Omega}\left(\sum_{j=1}^d\left(\frac{\partial u}{\partial x_j}\right)^2 + u^2\right)dx \leq C\int_{\Omega}f^2(x)\,dx.
\]

\item[b)] Let $\{u_n\}$ be a sequence of harmonic functions on $\Omega$, such that $\|u_n\|_{L^2(\Omega)}\leq M < \infty$ for a constant $M$ independent of $n$. Prove that there is a converging subsequence $\{u_{n_k}\}$ in $L^2(\Omega)$.
\end{itemize}
\end{exercise}

\begin{solution}
\textbf{Prerequisites:}
\par\medskip
1. \textbf{Green's First Identity:} This identity is a higher-dimensional version of integration by parts: $\int_\Omega (u \Delta v + \nabla u \cdot \nabla v) dx = \int_{\partial \Omega} u \frac{\partial v}{\partial n} ds$. For a function $u$ that vanishes on the boundary, this simplifies to $\int_\Omega u \Delta u dx = -\int_\Omega |\nabla u|^2 dx$.
\par\medskip
2. \textbf{Poincaré Inequality:} For a bounded domain $\Omega$, there exists a constant $C_P$ such that for all $u \in H_0^1(\Omega)$, $\int_\Omega u^2 dx \leq C_P \int_\Omega |\nabla u|^2 dx$. This theorem ensures that if a function's oscillations (gradient) are controlled, its magnitude ($L^2$ norm) is also controlled, provided it is "anchored" at the boundary.
\par\medskip
3. \textbf{Mean Value Property for Harmonic Functions:} If $\Delta u = 0$, then $u(x)$ is equal to the average of $u$ over any ball centered at $x$ contained in $\Omega$. A crucial consequence is that harmonic functions are smooth ($C^\infty$) and their derivatives at any point are bounded by their $L^2$ norm in a slightly larger neighborhood.
\par\medskip
4. \textbf{Arzelà-Ascoli Theorem:} This theorem states that a sequence of functions that is pointwise bounded and equicontinuous on a compact set has a uniformly convergent subsequence. For harmonic functions, $L^2$ boundedness implies local uniform boundedness and equicontinuity.
\par\medskip
\textbf{Steps:}

\textit{Part (a): Energy Estimates for the Poisson Equation}
\begin{enumerate}
    \item \textbf{Multiply and Integrate:} We start with the PDE $\Delta u = f$. Multiplying both sides by $u$ and integrating over the domain $\Omega$ gives:
    \[ \int_\Omega u \Delta u \, dx = \int_\Omega u f \, dx \]

    \item \textbf{Apply Integration by Parts:} Using Green's First Identity and the boundary condition $u|_{\partial \Omega} = 0$, the left side becomes:
    \[ -\int_\Omega |\nabla u|^2 \, dx = \int_\Omega u f \, dx \implies \int_\Omega |\nabla u|^2 \, dx = -\int_\Omega u f \, dx \]

    \item \textbf{Bound the Right-Hand Side:} Using the Cauchy-Schwarz inequality and then Young's inequality ($ab \leq \frac{\epsilon}{2} a^2 + \frac{1}{2\epsilon} b^2$), we have:
    \[ \int_\Omega |\nabla u|^2 \, dx \leq \int_\Omega |u||f| \, dx \leq \frac{\epsilon}{2} \int_\Omega u^2 \, dx + \frac{1}{2\epsilon} \int_\Omega f^2 \, dx \]

    \item \textbf{Invoke Poincaré Inequality:} Since $u \in H_0^1(\Omega)$, we substitute $\int u^2 \leq C_P \int |\nabla u|^2$:
    \[ \int_\Omega |\nabla u|^2 \, dx \leq \frac{\epsilon C_P}{2} \int_\Omega |\nabla u|^2 \, dx + \frac{1}{2\epsilon} \int_\Omega f^2 \, dx \]
    By choosing $\epsilon = 1/C_P$, we get $\frac{1}{2} \int |\nabla u|^2 \leq \frac{C_P}{2} \int f^2$, so $\int |\nabla u|^2 \leq C_P \int f^2$.

    \item \textbf{Sum the Norms:} Using Poincaré again, $\int u^2 \leq C_P \int |\nabla u|^2 \leq C_P^2 \int f^2$. Thus:
    \[ \int_\Omega (|\nabla u|^2 + u^2) dx \leq (C_P + C_P^2) \int_\Omega f^2 dx \]
    This proves the inequality with $C = C_P + C_P^2$.
\end{enumerate}

\textit{Part (b): Compactness of the Space of Harmonic Functions}
\textbf{The printed statement is false as an $L^2(\Omega)$ compactness claim.}
Take \(\Omega\) to be the unit disk in \(\mathbb R^2\), and define
\[
u_n(r,\theta)=\sqrt{\frac{2n+2}{\pi}}\,r^n\cos(n\theta).
\]
Each \(u_n\) is harmonic. Also,
\[
\|u_n\|_{L^2(\Omega)}^2
=\frac{2n+2}{\pi}\int_0^1r^{2n}r\,dr
\int_0^{2\pi}\cos^2(n\theta)\,d\theta
=1.
\]
For \(m\neq n\), the angular orthogonality of \(\cos(m\theta)\) and \(\cos(n\theta)\) gives
\[
\langle u_m,u_n\rangle_{L^2(\Omega)}=0.
\]
Thus \(\{u_n\}\) is an orthonormal sequence in \(L^2(\Omega)\). It has no strongly convergent subsequence, because
\[
\|u_m-u_n\|_{L^2}^2=2
\]
for \(m\neq n\).

\textbf{Correct compact local version.}
The usual true statement is local compactness: for every compact \(K\Subset\Omega\), a sequence of harmonic functions bounded in \(L^2(\Omega)\) has a subsequence converging uniformly, hence in \(L^2(K)\), on \(K\). This follows from the mean value property and interior gradient estimates, which give local uniform boundedness and equicontinuity; Arzel\`a--Ascoli and a diagonal argument then apply. The failure above comes from boundary concentration, so local compactness does not imply compactness in \(L^2(\Omega)\).
\par\medskip
\textbf{Intuition:}
\par\medskip
Part (a) is about the stability of physical systems. If you have a system (like heat distribution or a membrane) described by $\Delta u = f$, the "energy" required to maintain the state $u$ is bounded by the strength of the "force" $f$. It says you can't get an infinitely large response from a finite source.
\par\medskip
Part (b), however, shows the boundary limitation of that intuition. Harmonic functions are rigid in the interior, but on the full domain they can concentrate higher and higher oscillation near the boundary while keeping their \(L^2\) norm fixed.
\end{solution}
